\chapter{Conclusiones}
\label{cap:conclusiones}

En los veinte pares nocturnos del TNO, el algoritmo propuesto produce imágenes de calidad comparable a
la de los cinco métodos multiescala. Frente al Top-Hat clásico, la propuesta es mejor en cuatro de las
seis métricas. Lo que la distingue es su perfil. Por su puesto medio en cada grupo de métricas, queda primera en las de actividad, que miden
información, contraste y nitidez, y sexta de siete en las de fidelidad, que miden el parecido con las
fuentes. Con las seis métricas encabeza la clasificación por rango promedio, es decir, por el puesto
medio de cada método en todos los pares y métricas. Lo hace, sin embargo, por un margen estrecho y sin
una prueba que lo respalde. Ese primer lugar depende, además, del conjunto de métricas, y se obtuvo con
los métodos comparados en sus parámetros por defecto. Por eso no se concluye que la propuesta supere a los métodos multiescala,
sino que se ubica entre ellos con un perfil propio, que privilegia el realce sobre la fidelidad. En la
aplicación a la detección, el detector entrenado con las imágenes de la propuesta aventaja por un margen
mínimo al entrenado con las del Top-Hat clásico. Esa diferencia depende de la semilla aleatoria del
entrenamiento y queda como un hallazgo por confirmar.

Las conclusiones particulares siguen el orden de los objetivos específicos
(Sección~\ref{sec:objetivos}):

\begin{enumerate}
\item \textbf{Implementación y ajuste del algoritmo:} El algoritmo quedó implementado como se describe
en la Sección~\ref{sec:propuesta}, con el banco de cinco elementos estructurantes y la regla de fusión
con peso. Su radio y su peso se ajustaron mediante optimización por enjambre de partículas, con un
rango de búsqueda de 1 a 25 para el radio y de 0.30 a 2.00 para el peso. Las veinticinco corridas terminaron
con el peso en el límite inferior, $m = 0.30$. En el radio, en cambio, se repartieron sobre todo entre
los dos extremos del rango: dieciséis terminaron en $r = 1$, ocho en $r = 25$ y una en $r = 14$. Ese
valor del peso es un borde del rango y no un máximo interior: la grilla de comprobación mostró que
la aptitud, la función que el enjambre busca maximizar, era mayor con pesos menores que 0.30. El mejor
peso de cada radio estuvo entre 0.05 y 0.07. Se adoptó $r = 25$, uno de los dos resultados del
enjambre, porque prioriza el realce. Frente a $r = 1$, sobre los veinte pares y con las nueve métricas
iniciales, ese radio fue mejor en las cinco de actividad y peor en las cuatro de fidelidad. Esa elección
se hizo sobre las mismas imágenes con que después se evaluó la calidad, y explica en parte la ventaja
de la propuesta en actividad.

\item \textbf{Calidad de imagen:} Frente al Top-Hat clásico, la propuesta ganó en entropía, desviación
estándar, información mutua y similitud estructural. Perdió en gradiente medio y empató en la relación
señal-ruido de pico. Frente a los cinco métodos multiescala, su calidad fue comparable. Obtuvo el mejor
rango promedio, 3.642 frente a 3.783 de la DTCWT, pero la diferencia es de 0.14 puestos y no tiene
prueba propia. Con diecisiete métricas, en cambio, la propuesta quedó tercera
(Sección~\ref{sec:res-robustez}). Los contrastes de Wilcoxon frente a los métodos multiescala muestran
el mismo perfil: la propuesta ganó catorce de los quince de actividad y perdió trece de los quince de
fidelidad. Ese
intercambio entre realce y fidelidad proviene de la regla de fusión. En el barrido del peso, al subir
$m$ crecieron la desviación estándar y el gradiente medio y bajaron las tres métricas de fidelidad. Con
el Top-Hat clásico ocurrió lo mismo. La cabeza de la clasificación se mantuvo al tomar un bloque por escena. Con esos trece
bloques, la propuesta siguió primera y sesenta y cinco de los sesenta y seis contrastes conservaron su
decisión. El control con
versiones degradadas del promedio simple, en cambio, señala un límite de las métricas de actividad. Con
el ruido más intenso, ese promedio obtuvo el mejor valor de las catorce entradas en esas tres métricas.
Por eso, una ventaja en actividad no prueba, por sí sola, que el detalle agregado sea útil.

\item \textbf{Aplicación a la detección:} En el conjunto M3FD, el detector entrenado con las imágenes
de la propuesta alcanzó una tasa de acierto equilibrada del 61.4~\%. Esa tasa es el promedio entre el
acierto sobre los objetos que solo ve la cámara visible y el acierto sobre los que solo ve la infrarroja. La propuesta quedó
cuarta entre las ocho entradas fusionadas; la mayor tasa correspondió a la pirámide de Laplace, con
66.5~\% (Tabla~\ref{tab:v2_tasas}). Frente al Top-Hat clásico, la tasa fue 3.0 puntos mayor, con un
intervalo de confianza del 95~\% de 0.0 a 6.0 (Tabla~\ref{tab:deteccion_tophat}). Es una ventaja
mínima: el límite inferior, de 0.02 puntos, apenas supera el cero. Con los
otros dos umbrales de confianza también lo supera, pero no con el \textit{bootstrap} salvaje por grupos, otra forma de
calcular el intervalo, más confiable cuando los grupos son pocos y desiguales. Con el umbral principal, ese
cálculo dio un intervalo de $-3.5$ a 5.5 y un valor de $p$ de 0.17. La diferencia tampoco es robusta a la
semilla: según el par de semillas, fue de $-3.1$ a 7.8 puntos. Se trata, por lo tanto, de un hallazgo por confirmar con más grupos, más semillas y otros
detectores (Capítulo~\ref{cap:recomendaciones}). Aunque se confirmara, la ventaja no podría atribuirse
al banco, porque entre los dos métodos cambian a la vez los elementos estructurantes, el radio y el
peso. Tampoco se mostró que el detalle agregado ayude a detectar: frente al promedio simple, la
diferencia fue de 1.2 puntos, con un intervalo de $-6.6$ a 9.0.

\item \textbf{Examen de la hipótesis:} En su parte de calidad, la hipótesis planteada en el
Capítulo~\ref{cap:introduccion} se sostiene en lo principal (Sección~\ref{sec:res-hipotesis}). La propuesta superó
al Top-Hat clásico en cuatro de las seis métricas, y su calidad resultó comparable a la de los cinco
métodos multiescala. En la aplicación a la detección, la ventaja frente al Top-Hat clásico es mínima.
Con el umbral principal, el límite inferior del intervalo supera el cero por muy poco, y la diferencia
cambia de signo según el par de semillas. Esa parte de la hipótesis queda, por eso, como un hallazgo
que falta confirmar.
\end{enumerate}

El aporte del trabajo es un algoritmo de fusión de imágenes infrarrojas y visibles que puede explicarse
paso a paso. El algoritmo lleva a esta fusión el banco de cinco elementos estructurantes que
\textcite{floressf} emplearon para realzar mamografías, y combina sus respuestas según
\textcite{bala2024}. La regla con peso y el ajuste del radio y del peso por enjambre de partículas
siguen a \textcite{ortegarodriguez2025}. Cada parámetro tiene una lectura directa: el radio determina el
tamaño del detalle que se extrae, y el peso, cuánto de ese detalle se suma. La calidad de sus imágenes
quedó caracterizada con seis métricas y pruebas no paramétricas. Frente a los métodos multiescala, esas
imágenes ganan contraste y nitidez y pierden fidelidad a las fuentes, en una medida que regula el peso.
La aplicación a la detección se evaluó con detectores entrenados sobre cada tipo de imagen, y la ventaja
frente al Top-Hat clásico quedó por confirmar. No se halló una relación significativa entre el puesto
de cada método en calidad y su tasa equilibrada (Sección~\ref{sec:disc-relacion}). Los dos métodos
morfológicos, los mejores en actividad, quedaron cuarto y séptimo de ocho en detección. Esa
falta de correspondencia respalda medir la fusión también por su efecto en una tarea que use sus
imágenes \parencite{liu2022, tang2022}.
